\documentclass[12pt]{article}

\usepackage{lipsum}
\usepackage[margin=1in]{geometry}
\usepackage{graphicx}
\graphicspath{{../data}}
\usepackage{float} % for H specifier
\usepackage{subcaption} % subfigures

\usepackage{xcolor}

%\usepackage{endfloat} 
%nomarkers
%nolists
%figuresonly
%\renewcommand{\efloatseparator}{\bigskip}

\begin{document}

\section{Introduction}
Images in \LaTeX{} are handled by the \texttt{graphicx} package. We include the actual image file using \texttt{\textbackslash includegraphics}.

However, we usually wrap the image in a \texttt{\textbackslash begin\{figure\}} environment. This is called a \textbf{float}. \texttt{\textbackslash begin\{table\}} is another example of a float.  
Floats drift to the optimal location (top of page, bottom, etc.)to prevent ugly page breaks. You guide this behavior using \textbf{placement specifiers}:
\begin{itemize}
    \item \texttt{h} (here) - Place around here (respecting `aesthetics').
    \item \texttt{t} (top), \texttt{b} (bottom), \texttt{p} (page of floats).
    \item \texttt{H} (HERE) - Exactly here, (requires the \texttt{float} package).
\end{itemize}

\section{Single Images}
\lipsum[1-5]
\begin{figure} % change
    \centering
    \includegraphics[width=0.8\linewidth]{images/greenCells.png}
    \caption{This is an image of  green cells. There is another figure (Figure \ref{fig:grayCells}) full of gray cells! }
    \label{fig:greenCells}
\end{figure}

{\color{red}
\lipsum[4-9]
}
%!h -> forget aesthetic put around here
% H -> put here, dont care about nothin
% h -> put around here respecting aesthetics
\begin{figure}[!t]%[H]%[!h]%[!htbp] %[h] v/s [!h] v/s [H] push all to end coz figure voilates limits v/s
    \centering
    \includegraphics[width=\textwidth]{images/blue.jpg}
    \caption{The blue figure}
    \label{fig:blueCross}
\end{figure}

{\color{green}
\lipsum[1-5]
}
\begin{figure}[!htbp]
    \centering
    \includegraphics[width=0.6\textwidth]{images/yellow.jpg}
    \caption{The yellow figure}
    \label{fig:yellowCross}
\end{figure}

{\color{blue}
\lipsum[1-5]
}
\begin{figure}[!htbp]
    \centering
    \includegraphics[width=0.6\textwidth]{images/red.jpg}
    \caption{The red figure}
    \label{fig:redCross}
\end{figure}

{\color{orange}
\lipsum[1-5]
}
\begin{figure}[!htbp]
    \centering
    \includegraphics[width=0.6\textwidth]{images/cyan.jpg}
    \caption{The cyan figure}
    \label{fig:cyanCross}
\end{figure}

{\color{purple}
\lipsum[1-5]
}

\begin{figure}[!htbp]
    \centering
    \includegraphics[width=0.6\textwidth]{images/orange.jpg}
    \caption{The orange figure}
    \label{fig:orangeCross}
\end{figure}

{\color{violet}
\lipsum[1-5]
}

\begin{figure}[htbp]
    \centering
    \includegraphics[width=\linewidth]{images/grayCells.png}
    \caption{This is an image of gray cells pretty drab compared to the colorfull Figure \ref{fig:greenCells}}
    \label{fig:grayCells}
\end{figure}

{\color{green!50!blue}
\lipsum[10-17]
}

\twocolumn
\section{Image in two column documents}
\lipsum[1]
\begin{figure}[h]
    \centering
    \includegraphics[width=0.5\linewidth]{images/blue.jpg}
    \caption{Blue Cross}
    \label{fig:2colBlueCross}
\end{figure}
\color{red}\lipsum[3-9]
\begin{figure*}[h]
    \centering
    \includegraphics[width=\linewidth]{images/grayCells.png}
    \caption{Across the two column}
    \label{fig:acrossColumns}
\end{figure*}
\color{blue}\lipsum[1-7]
\end{document}
